\slajdid{09_3-s061}
\begin{frame}[label=09_3-s061]
\gusttekst\setlength{\abovedisplayskip}{3pt}\setlength{\belowdisplayskip}{3pt}\setlength{\abovedisplayshortskip}{2pt}\setlength{\belowdisplayshortskip}{2pt}Slučaj $V_C=+15V,V_I=-12V$\begin{center}\input{slike/tikz/s061_bilateralni_prekidac.tex}\end{center}N tranzistor sigurno ima uslove za provođenje tako što će mu sors biti na strani 1. Radi sa velikim naponima između gejta i sorsa i sa velikom strujom. Struja kroz $R_L$ će teći u smeru prikazanom na slici pa je napon na drejnu N tranzistora negativan i on sigurno radi u omskoj oblasti. Ako je otpornost $R_L$ mnogo veća od dinamičke otpornosti N tranzistora koji radi u omskoj oblasti $V_O\approx-12V$, bez degradacije naponskih nivoa. Bitno je da uočite šta se dešava sa P tranzistorom. Pod pretpostavkom da je $V_{Tn}\approx-V_{Tp}<3V$ i tranzistor P ima uslove za provođenje. Njegov napon između gejta i sorsa je $V_{GS}<V_{Tp}$. Sors mu je na strani 2. Zbog malog napona između drejna i sorsa (obezbeđuje N tranzistor) on je takođe u omskoj oblasti, ali zbog relativno malog napona između sorsa i gejta sa velikom otpornošću.
\end{frame}
